% Auto-generated complete chapter from 13_interview_negotiation_and_offer_risk.md
\chapter{13 — Interview, Negotiation and Offer-Risk Framework}
\label{complete:root_015}

\begin{sourcemeta}
\textbf{Fuente:} \href{file:///E:/Laboral/13_interview_negotiation_and_offer_risk.md}{13\_interview\_negotiation\_and\_offer\_risk.md}\\
\textbf{Seccion:} root\\
\textbf{Estado:} canonical\\
\textbf{Rol:} Modulo principal\\
\textbf{Lineas originales:} 816\\
\textbf{Ultima modificacion:} 2026-06-11 14:40\\
\textbf{Deduplicacion conservadora:} 0 bloques repetidos omitidos; 0 caracteres repetidos no reimpresos.
\end{sourcemeta}

\section*{Resumen de orientacion}
This file defines a practical interview, negotiation and offer-risk framework for a Colombia-based Unity / Real-Time 3D / Technical Visualization candidate.

\section*{Contenido completo depurado}
\addcontentsline{toc}{section}{Contenido completo depurado}




\subsection{1. Purpose and scope}




This file defines a practical interview, negotiation and offer-risk framework for a Colombia-based Unity / Real-Time 3D / Technical Visualization candidate.




It is built around the source-of-truth profile established in previous files:




\begin{itemize}
\item Flagship project: \textbf{TwinSight X500}, a Unity WebGL technical-visualization thesis project.
\item Positioning: Real-Time 3D / Unity Technical Artist / Technical Visualization / WebGL / CAD-to-realtime.
\item Current legal position: Colombia-based remote contractor or employee through compliant hiring structure.
\item Future mobility: Portuguese passport expected in approximately two years; Germany route possible through marriage and residence processing.
\item Experience profile: strong portfolio evidence, limited formal industry experience.
\item Languages: Spanish native; English C1 self-assessed; German not currently usable professionally.
\end{itemize}




This document is not legal, tax or immigration advice. It is an operational framework for interviews, negotiation, contract review and offer evaluation.




\subsection{2. Source-quality audit of the Deep Research output}




The Deep Research output was useful as a raw input, but not acceptable as a final section without correction.




\subsubsection{What was useful}




\begin{itemize}
\item It identified common interview formats: screening, technical interview, take-home test, hiring-manager interview and offer stage.
\item It surfaced representative Unity interview topics: GameObjects, components, \texttt{Update}, \texttt{FixedUpdate}, prefabs and colliders.
\item It correctly separated employee and contractor negotiation issues.
\item It highlighted red flags in remote offers.
\item It connected salary anchoring to section 03 benchmarks.
\end{itemize}




\subsubsection{What required correction}




\begin{itemize}
\item It repeated salary data that belongs primarily in section 03.
\item It used some weak sources for interview patterns and red flags.
\item It overstated some tax points, especially around IVA/exported services, which must be handled cautiously.
\item It did not adequately separate role-family interview expectations.
\item It did not provide a rigorous take-home test risk model.
\item It did not give a structured offer-scoring model.
\item It did not sufficiently address AI-assisted development disclosure.
\end{itemize}




\subsubsection{Corrective direction}




This final version keeps the useful structure, removes weak generalizations, and organizes the section around the actual job-search process:




\begin{enumerate}
\item Interview evidence package.
\item Interview preparation by role family.
\item Technical proof expectations.
\item Take-home and unpaid-labor risk.
\item Negotiation logic.
\item Contractor and compliance risks.
\item AI-disclosure handling.
\item Offer evaluation.
\item Accept / counter / reject rules.
\end{enumerate}




\subsection{3. Candidate-specific interview positioning}




The candidate should not enter interviews as a generic Unity developer, generic 3D artist or generalist full-stack developer.




The most defensible interview frame is:




\begin{quote}
Real-Time 3D Developer / Unity Technical Artist focused on interactive technical visualization, CAD-to-realtime optimization, WebGL deployment and technical UI systems.
\end{quote}




\subsubsection{Defensible claims}




\begin{itemize}
\item Built a Unity WebGL interactive 3D technical visualization prototype for drone assembly understanding.
\item Worked across CAD-to-realtime optimization, Unity integration, interaction systems, UI, visual modes and technical documentation.
\item Used academic usability and workload metrics to evaluate the tool.
\item Can explain asset optimization trade-offs, interaction design, WebGL constraints and validation methodology.
\item Used AI as an implementation assistant while retaining technical judgment, debugging, integration and final ownership.
\end{itemize}




\subsubsection{Claims to avoid}




\begin{itemize}
\item Do not claim senior industry experience.
\item Do not claim shipped commercial product experience unless a client/product actually shipped.
\item Do not claim current EU work authorization.
\item Do not claim pure AI engineer / ML engineer profile.
\item Do not claim Houdini FX TD ability unless a public reel proves it.
\item Do not imply that all TwinSight code was handcrafted line-by-line without AI assistance if challenged directly.
\end{itemize}




\subsection{4. Interview evidence package}




Before applying or interviewing, the candidate should have a compact evidence package.




\subsubsection{Minimum evidence package}




\begin{center}
\scriptsize
\begin{longtable}{>{\raggedright\arraybackslash}p{0.455\linewidth} >{\raggedright\arraybackslash}p{0.455\linewidth}}
\toprule
Asset & Required state \\
\midrule
\endfirsthead
\toprule
Asset & Required state \\
\midrule
\endhead
Portfolio page & TwinSight case study public \\
Demo video & 60–120 seconds \\
GitHub README & Recruiter-readable \\
Screenshots & Clear UI and features \\
Metrics table & Triangles, SUS, NASA-TLX \\
Technical diagram & CAD-to-realtime pipeline \\
Role statement & Exact contribution stated \\
AI disclosure & Prepared but not overemphasized \\
\bottomrule
\end{longtable}
\end{center}




\subsubsection{Strong evidence package}




The stronger version adds:




\begin{itemize}
\item WebGL live build or controlled video demo.
\item Before/after optimization visuals.
\item Interaction-system architecture diagram.
\item WebGL constraints section.
\item Mobile-performance discussion.
\item One clean C\# code sample.
\item One editor-tool or automation sample.
\item One concise technical breakdown on ArtStation or portfolio site.
\end{itemize}




\subsection{5. Interview pipeline by role family}




\subsubsection{5.1 Unity Technical Artist}




Likely interview focus:




\begin{itemize}
\item Unity rendering pipeline: URP, materials, Shader Graph, basic HLSL.
\item Asset optimization: meshes, materials, textures, batching, draw calls, LODs.
\item Art-to-engine integration.
\item Debugging visual bugs.
\item Technical communication between artists and programmers.
\item Tooling or editor scripts.
\end{itemize}




Expected proof:




\begin{itemize}
\item Visual breakdown of TwinSight modes.
\item CAD-to-realtime before/after examples.
\item Explanation of optimization decisions.
\item One shader or visual-mode breakdown.
\item One example of bridging asset pipeline and runtime constraints.
\end{itemize}




Risk:




The profile is viable for optimization/runtime/integration technical art, but weaker for shader-heavy, Houdini-heavy or AAA pipeline roles.




\subsubsection{5.2 Real-Time 3D Developer}




Likely interview focus:




\begin{itemize}
\item C\# architecture.
\item Runtime interaction systems.
\item Object selection, highlighting, isolation, camera controls.
\item UI integration.
\item Data-driven component panels.
\item Performance profiling.
\item WebGL deployment constraints.
\end{itemize}




Expected proof:




\begin{itemize}
\item TwinSight interaction system explanation.
\item Technical UI architecture.
\item Build constraints and optimization notes.
\item Clear answer to “what did you program yourself?”
\end{itemize}




Risk:




The candidate must avoid sounding like a pure artist. The value proposition is technical implementation, interaction design and 3D systems.




\subsubsection{5.3 Unity WebGL / Web 3D Developer}




Likely interview focus:




\begin{itemize}
\item Unity WebGL build constraints.
\item Memory limits and load time.
\item Browser compatibility.
\item Asset compression.
\item Asynchronous loading.
\item Deployment workflows.
\item Front-end integration if the role involves React/Three.js.
\end{itemize}




Expected proof:




\begin{itemize}
\item WebGL build or video evidence.
\item Explanation of WebGL limitations encountered.
\item Asset-size and performance trade-offs.
\item Understanding of mobile browser constraints.
\end{itemize}




Risk:




If the role is actually Three.js/WebGPU-heavy, TwinSight helps conceptually but does not fully prove the stack. The candidate should present Three.js/WebGPU as an adjacent skill only if supported by a small demo.




\subsubsection{5.4 Technical Visualization / Digital Twin / Simulation Developer}




Likely interview focus:




\begin{itemize}
\item Translating physical systems into interactive 3D representations.
\item CAD-to-realtime workflow.
\item Scene hierarchy and component relationships.
\item Technical UI and metadata.
\item Simulation vs visualization distinction.
\item Data integration concepts.
\item Usability for technical users.
\end{itemize}




Expected proof:




\begin{itemize}
\item TwinSight as the primary case study.
\item Explanation of how 2D technical documentation becomes interactive 3D.
\item Metrics and user-evaluation results.
\item Clear limitation statement: TwinSight is not a physics-accurate digital twin unless real-time sensor/data integration exists.
\end{itemize}




Risk:




Do not overclaim “digital twin” if the system does not connect live operational data. Use “technical visualization” or “digital-twin-adjacent prototype” when necessary.




\subsubsection{5.5 XR / AR / VR Developer}




Likely interview focus:




\begin{itemize}
\item Unity XR stack.
\item Spatial interaction.
\item Performance constraints.
\item Input systems.
\item Comfort and usability.
\item Platform SDKs.
\end{itemize}




Expected proof:




\begin{itemize}
\item TwinSight proves 3D interaction and technical visualization, not XR experience by itself.
\item A small XR prototype would significantly improve credibility.
\item If no XR artifact exists, position XR as adjacent, not primary.
\end{itemize}




Risk:




XR roles may require headset-specific experience. Without Meta Quest / OpenXR / AR Foundation artifacts, probability is lower.




\subsubsection{5.6 Tools / Pipeline Developer}




Likely interview focus:




\begin{itemize}
\item Automation mindset.
\item Python scripting.
\item Unity editor tooling.
\item Asset processing.
\item Naming, validation, import rules.
\item Batch operations.
\item Pipeline documentation.
\end{itemize}




Expected proof:




\begin{itemize}
\item ARA as secondary tooling evidence.
\item Any Unity editor tools or scripts used in TwinSight.
\item Before/after workflow reduction.
\item Small public tool repo if possible.
\end{itemize}




Risk:




ARA supports tooling logic, but does not prove production art-pipeline engineering unless made stable, documented and demoable.




\subsubsection{5.7 Python Automation / AI Tools / LLM Application Developer}




Likely interview focus:




\begin{itemize}
\item Python architecture.
\item API use.
\item LangGraph or LLM orchestration.
\item Error handling.
\item Prompt/tool reliability.
\item Data flow.
\item Cost control.
\item Evaluation.
\end{itemize}




Expected proof:




\begin{itemize}
\item ARA demo or README.
\item Clear architecture diagram.
\item Known limitations.
\item No overclaiming as ML engineer.
\end{itemize}




Risk:




This route can dilute the main profile. It should be used for tools, research automation and technical workflow roles, not as a full AI engineer identity.




\subsection{6. Technical question bank}




This section is not intended for memorization. It defines categories that should be practiced with precise answers.




\subsubsection{Unity fundamentals}




\begin{itemize}
\item What is a GameObject and how do components relate to it?
\item Explain the Unity lifecycle: \texttt{Awake}, \texttt{Start}, \texttt{Update}, \texttt{FixedUpdate}, \texttt{LateUpdate}.
\item When should \texttt{FixedUpdate} be used?
\item What is a prefab and why does it matter in production?
\item What are colliders and rigidbodies?
\item How does Unity serialization work at a basic level?
\item How do ScriptableObjects differ from MonoBehaviours?
\item How would you structure object selection and highlighting?
\end{itemize}




\subsubsection{Performance and optimization}




\begin{itemize}
\item How do you diagnose low FPS in Unity?
\item What metrics do you check first: CPU, GPU, draw calls, batches, memory, textures, triangles?
\item How would you optimize a CAD-imported scene?
\item What is the difference between static batching, dynamic batching, GPU instancing and SRP batching?
\item When should you use LODs?
\item How do texture size and material count affect WebGL performance?
\item What causes long load times in a browser build?
\item How do you reduce memory pressure in Unity WebGL?
\end{itemize}




\subsubsection{Rendering and technical art}




\begin{itemize}
\item What is URP and why use it?
\item What can Shader Graph solve without custom HLSL?
\item When is HLSL necessary?
\item How would you implement X-Ray, ghosted, wireframe or blueprint visual modes?
\item What are common shader performance problems?
\item What makes a material setup maintainable?
\end{itemize}




\subsubsection{Technical visualization}




\begin{itemize}
\item How do you transform technical documentation into an interactive 3D experience?
\item What is the difference between product visualization, technical visualization, simulation and a digital twin?
\item How would you organize part metadata?
\item How would you design a component hierarchy for a complex assembly?
\item How would you present uncertainty or non-physical visualization modes such as qualitative thermal overlays?
\end{itemize}




\subsubsection{WebGL / web deployment}




\begin{itemize}
\item What are Unity WebGL constraints?
\item How do browsers affect memory and performance?
\item How would you optimize first-load experience?
\item What is the difference between Unity WebGL and Three.js for web-based 3D?
\item What should be tested on mobile browsers?
\end{itemize}




\subsubsection{Tools / pipeline}




\begin{itemize}
\item What manual process would you automate first in a 3D pipeline?
\item How would you validate imported assets?
\item How would you batch-process hundreds of meshes or textures?
\item How would you design a naming convention or asset audit tool?
\item How would you document a pipeline for artists and engineers?
\end{itemize}




\subsubsection{AI tooling}




\begin{itemize}
\item What problem does ARA solve?
\item What are the failure modes of LLM-based automation?
\item How do you evaluate the output of an AI tool?
\item How do you avoid hallucinated citations or unverifiable claims?
\item How do you control cost and retries in LLM workflows?
\end{itemize}




\subsection{7. TwinSight interview narrative}




\subsubsection{60-second version}




\begin{quote}
TwinSight X500 is my thesis project: a Unity WebGL interactive 3D technical-visualization prototype for understanding the assembly of a Holybro X500 V2 drone. The goal was to move beyond static 2D documentation and help users understand spatial relationships between components. I worked on the Unity integration, interaction systems, UI, visual modes, CAD-to-realtime optimization and evaluation documentation. The project reduced heavy CAD geometry into a realtime-ready asset and was evaluated with usability and workload instruments.
\end{quote}




\subsubsection{2-minute version}




The extended explanation should cover:




\begin{enumerate}
\item Problem: 2D documentation requires mental reconstruction.
\item Solution: interactive 3D inspection and assembly understanding.
\item Technical stack: Unity, WebGL, C\#, URP, UI Toolkit, Blender pipeline.
\item Pipeline: CAD/STEP to optimized real-time asset.
\item Features: selection, isolation, exploded view, clipping, visual modes, part data.
\item Evaluation: SUS, NASA-TLX Raw, Think-Aloud.
\item Contribution: integration, systems, pipeline, documentation, debugging, AI-assisted implementation.
\item Limits: prototype, not shipped commercial product; not live IoT digital twin.
\end{enumerate}




\subsubsection{Evidence statements}




Use quantified claims only when they are in the thesis or portfolio material:




\begin{itemize}
\item CAD route above 6.5 million triangles reduced to 95,617 optimized triangles.
\item 12 validation participants.
\item SUS average: 91.88.
\item NASA-TLX Raw: 8.69 for 3D viewer vs 19.89 for 2D support.
\item 96 task-condition records.
\end{itemize}




\subsection{8. Take-home test risk analysis}




Take-home tests can be useful, but they carry unpaid-labor and scope-risk problems.




\subsubsection{Acceptable take-home test}




\begin{center}
\scriptsize
\begin{longtable}{>{\raggedright\arraybackslash}p{0.455\linewidth} >{\raggedright\arraybackslash}p{0.455\linewidth}}
\toprule
Criterion & Acceptable form \\
\midrule
\endfirsthead
\toprule
Criterion & Acceptable form \\
\midrule
\endhead
Scope & 2–4 hours maximum \\
Objective & Skill demonstration \\
Dataset & Toy or anonymized data \\
IP & Candidate retains portfolio rights if generic \\
Feedback & Some review or next step \\
Deadline & Reasonable, not overnight \\
\bottomrule
\end{longtable}
\end{center}




\subsubsection{Risky take-home test}




\begin{center}
\scriptsize
\begin{longtable}{>{\raggedright\arraybackslash}p{0.455\linewidth} >{\raggedright\arraybackslash}p{0.455\linewidth}}
\toprule
Risk & Signal \\
\midrule
\endfirsthead
\toprule
Risk & Signal \\
\midrule
\endhead
Free labor & Production-ready feature requested \\
Unbounded scope & “Spend as much time as needed” \\
IP capture & Company keeps all code with no compensation \\
No feedback & Auto-rejection after work \\
Vague evaluation & No criteria \\
Excessive urgency & 24-hour turnaround for complex task \\
\bottomrule
\end{longtable}
\end{center}




\subsubsection{Refusal or boundary script}




Use a firm but neutral boundary:




\begin{quote}
I’m comfortable completing a focused technical assessment. To keep it fair, I can spend up to 3–4 hours and provide a short README explaining decisions and trade-offs. If the task requires production-ready implementation beyond that scope, I’d prefer to treat it as a paid trial or reduce the deliverable.
\end{quote}




\subsubsection{Paid-trial preference}




For complex real company work:




\begin{quote}
This looks close to a production task. I can do it as a paid trial with a defined scope, deliverables, IP terms and timeline.
\end{quote}




\subsection{9. Negotiation framework}




\subsubsection{9.1 Core negotiation principle}




For international remote work from Colombia, the candidate should negotiate based on \textbf{gross contractor compensation}, not local Colombian employee salary.




A contractor must cover:




\begin{itemize}
\item taxes;
\item social-security contributions;
\item health insurance or EPS obligations;
\item retirement/pension contributions;
\item unpaid vacation;
\item unpaid sick days;
\item equipment;
\item software subscriptions;
\item currency conversion and transfer fees;
\item accountant/legal support;
\item bench time between contracts.
\end{itemize}




Therefore, the contractor rate must be higher than an equivalent employee salary.




\subsubsection{9.2 Salary anchoring by market}




Use section 03 as the primary salary source. This file should not replace it.




\paragraph{Practical anchor bands}



\begin{center}
\scriptsize
\begin{longtable}{>{\raggedright\arraybackslash}p{0.305\linewidth} >{\raggedright\arraybackslash}p{0.305\linewidth} >{\raggedright\arraybackslash}p{0.305\linewidth}}
\toprule
Target & Gross monthly & Typical route \\
\midrule
\endfirsthead
\toprule
Target & Gross monthly & Typical route \\
\midrule
\endhead
Floor & USD 1.5k & local/LATAM junior role \\
Base target & USD 3k & LATAM contractor / junior-strong remote \\
Strong target & USD 4k–6k & mid-junior contractor / niche Unity-WebGL \\
High target & USD 6k+ & US/EU contractor, digital twin, simulation, XR \\
\bottomrule
\end{longtable}
\end{center}




\subsubsection{9.3 When asked “What are your salary expectations?”}




Do not give a single number too early if the scope is unclear.




Recommended answer:




\begin{quote}
It depends on scope, employment model and benefits. For a contractor arrangement, I usually evaluate the monthly gross range based on responsibilities, time commitment and whether the role is closer to Unity development, technical art, simulation or WebGL deployment. For a full-time contractor role, I would expect a range aligned with international remote Unity/real-time 3D roles rather than local Colombian payroll rates.
\end{quote}




If pressed:




\begin{quote}
For contractor work, my current target range is USD [X–Y] gross per month, depending on scope, exclusivity, payment terms and whether the work includes production ownership or exploratory prototyping.
\end{quote}




\subsubsection{9.4 Anchor ranges by priority}




\begin{center}
\scriptsize
\begin{longtable}{>{\raggedright\arraybackslash}p{0.305\linewidth} >{\raggedright\arraybackslash}p{0.305\linewidth} >{\raggedright\arraybackslash}p{0.305\linewidth}}
\toprule
Priority & Use case & Starting anchor \\
\midrule
\endfirsthead
\toprule
Priority & Use case & Starting anchor \\
\midrule
\endhead
A & US/EU technical visualization & USD 5k–7k+ \\
B & LATAM/global contractor & USD 3k–5k \\
C & local or fallback & USD 1.5k–3k \\
Paid trial & limited test & hourly/day rate \\
\bottomrule
\end{longtable}
\end{center}




These anchors must be adjusted after validating scope, seniority, hours and hiring model.




\subsection{10. Contractor vs employee decision matrix}




\begin{center}
\scriptsize
\begin{longtable}{>{\raggedright\arraybackslash}p{0.225\linewidth} >{\raggedright\arraybackslash}p{0.225\linewidth} >{\raggedright\arraybackslash}p{0.225\linewidth} >{\raggedright\arraybackslash}p{0.225\linewidth}}
\toprule
Model & Advantages & Risks & Best use \\
\midrule
\endfirsthead
\toprule
Model & Advantages & Risks & Best use \\
\midrule
\endhead
Colombia contractor & fastest setup & taxes, no benefits & foreign remote clients \\
EOR employee & compliant payroll & employer cost & serious foreign employer \\
Local employee & stability & lower pay & bridge role \\
Freelance platform & fast entry & low trust, fees & small projects \\
EU employee & benefits, EU rights & future only & after passport/residence \\
German family route & work authorization path & paperwork & if relocation occurs \\
\bottomrule
\end{longtable}
\end{center}




\subsubsection{Contractor checklist}




Before accepting a contractor role, verify:




\begin{itemize}
\item legal entity and signatory;
\item payment currency;
\item payment schedule;
\item late payment clause;
\item invoice requirements;
\item IP ownership;
\item confidentiality scope;
\item non-compete scope;
\item termination notice;
\item exclusivity;
\item equipment/software responsibility;
\item allowed public portfolio mention;
\item governing law;
\item dispute resolution;
\item tax documents required.
\end{itemize}




\subsection{11. Colombia-specific remote-contracting risks}




This section is high-level only. It must be confirmed with a Colombian accountant.




\subsubsection{Risk categories}




\begin{center}
\scriptsize
\begin{longtable}{>{\raggedright\arraybackslash}p{0.455\linewidth} >{\raggedright\arraybackslash}p{0.455\linewidth}}
\toprule
Risk & What to check \\
\midrule
\endfirsthead
\toprule
Risk & What to check \\
\midrule
\endhead
RUT & correct activity codes \\
Invoicing & DIAN electronic invoicing \\
IVA & export-service treatment \\
Retención & local vs foreign client \\
Social security & EPS, pension, ARL base \\
Currency & exchange and documentation \\
Tax residency & Colombian obligations \\
Misclassification & contractor vs employee behavior \\
\bottomrule
\end{longtable}
\end{center}




\subsubsection{Misclassification risk}




A contractor relationship becomes riskier if the company controls:




\begin{itemize}
\item fixed schedule;
\item daily supervision;
\item mandatory exclusivity;
\item company equipment only;
\item internal hierarchy;
\item long-term dependency;
\item employee-like benefits without formal employment.
\end{itemize}




If the role behaves like employment, the company may need an EOR or local employment structure.




\subsubsection{Practical mitigation}




Use a contract that states:




\begin{itemize}
\item independent contractor status;
\item deliverables or services;
\item invoicing procedure;
\item autonomy of execution;
\item no automatic employment benefits;
\item clear payment terms;
\item IP transfer only for paid deliverables.
\end{itemize}




\subsection{12. EU move implications on compensation}




\subsubsection{Current Colombia position}




Current strongest route:




\begin{itemize}
\item remote contractor;
\item foreign client;
\item USD or EUR gross compensation;
\item tax/social-security planning in Colombia.
\end{itemize}




\subsubsection{After Portuguese passport}




Once the passport is issued, the candidate can legally state EU work authorization, but compensation depends on location and employment model.




Possible outcomes:




\begin{itemize}
\item EU employee salary may be lower gross than US contractor compensation.
\item Benefits, healthcare and stability improve.
\item Cost of living and tax burden may increase.
\item EU work authorization improves probability for Germany, Portugal, Netherlands and EU-only remote roles.
\end{itemize}




\subsubsection{Germany through marriage/residence}




If legally married and resident in Germany under the correct residence title, Germany becomes a realistic work market.




Implications:




\begin{itemize}
\item more local/hybrid opportunities;
\item German language matters more;
\item health insurance and tax obligations become significant;
\item employer credibility improves because work authorization is clearer;
\item net income may be lower than gross headline salary suggests.
\end{itemize}




\subsubsection{Recruiter-facing wording by state}




\begin{center}
\scriptsize
\begin{longtable}{>{\raggedright\arraybackslash}p{0.455\linewidth} >{\raggedright\arraybackslash}p{0.455\linewidth}}
\toprule
State & Accurate wording \\
\midrule
\endfirsthead
\toprule
State & Accurate wording \\
\midrule
\endhead
Now & Colombia-based; available as remote contractor \\
Passport pending & EU citizenship expected in approximately two years, not current work authorization \\
Passport issued & EU citizen; authorized to work in EU/EEA \\
Germany residence & Authorized to work in Germany, if residence title grants it \\
Blue Card path & Would require sponsorship if used before EU citizenship \\
\bottomrule
\end{longtable}
\end{center}




\subsection{13. AI-assisted development disclosure}




AI use is not automatically disqualifying. Poor disclosure is the risk.




\subsubsection{Bad disclosure}




Avoid:




\begin{quote}
I used AI to build the project.
\end{quote}




This sounds like lack of ownership.




\subsubsection{Better disclosure}




Use:




\begin{quote}
I used AI as an implementation assistant for iteration, code generation support and debugging. I was responsible for architecture decisions, integration, testing, debugging, technical judgment and final behavior.
\end{quote}




\subsubsection{When asked “How much did you build yourself?”}




Answer structure:




\begin{enumerate}
\item State ownership.
\item Separate AI assistance from engineering responsibility.
\item Explain a concrete decision you made.
\item Mention debugging and validation.
\item Offer to walk through code or architecture.
\end{enumerate}




Example:




\begin{quote}
The project was AI-assisted, but I owned the system design, Unity integration, debugging, interaction behavior, asset workflow decisions and final validation. I used AI the way many developers use autocomplete, documentation search or pair-programming support: as an accelerator, not as a replacement for technical judgment. I can walk through the interaction architecture and explain the decisions behind the pipeline.
\end{quote}




\subsubsection{Evidence needed}




To make the disclosure credible, prepare:




\begin{itemize}
\item architecture diagram;
\item one code walkthrough;
\item one debugging story;
\item one performance trade-off;
\item one limitation you identified independently;
\item one feature you revised after testing.
\end{itemize}




\subsection{14. Offer red flags}




\subsubsection{Critical red flags}




\begin{center}
\scriptsize
\begin{longtable}{>{\raggedright\arraybackslash}p{0.455\linewidth} >{\raggedright\arraybackslash}p{0.455\linewidth}}
\toprule
Red flag & Risk \\
\midrule
\endfirsthead
\toprule
Red flag & Risk \\
\midrule
\endhead
No written offer & unenforceable terms \\
Salary changes after interview & bait-and-switch \\
unpaid production task & free labor \\
asks for money & scam \\
unclear company identity & fraud \\
vague IP terms & portfolio/legal risk \\
unlimited non-compete & career restriction \\
60+ day payment terms & cash-flow risk \\
crypto-only payment & compliance risk \\
“remote” but no country scope & work-authorization risk \\
\bottomrule
\end{longtable}
\end{center}




\subsubsection{Moderate red flags}




\begin{itemize}
\item no technical interviewer;
\item no clarity on reporting manager;
\item unrealistic timeline;
\item “we are like a family” language replacing contract clarity;
\item no clear definition of done;
\item no tool/license budget;
\item no test device or hardware plan for XR/WebGL;
\item “equity instead of pay” without legal documentation.
\end{itemize}




\subsection{15. Offer evaluation scorecard}




Use this scoring model before accepting any offer.




\subsubsection{Core scoring table}




\begin{center}
\scriptsize
\begin{longtable}{>{\raggedright\arraybackslash}p{0.305\linewidth} >{\raggedright\arraybackslash}p{0.305\linewidth} >{\raggedright\arraybackslash}p{0.305\linewidth}}
\toprule
Category & Weight & Score 1–5 \\
\midrule
\endfirsthead
\toprule
Category & Weight & Score 1–5 \\
\midrule
\endhead
Role fit & 20\% &  \\
Compensation & 20\% &  \\
Remote legality & 15\% &  \\
Portfolio leverage & 15\% &  \\
Contract safety & 15\% &  \\
Growth potential & 10\% &  \\
Lifestyle/time-zone fit & 5\% &  \\
\bottomrule
\end{longtable}
\end{center}




\subsubsection{Decision interpretation}




\begin{center}
\scriptsize
\begin{longtable}{>{\raggedright\arraybackslash}p{0.455\linewidth} >{\raggedright\arraybackslash}p{0.455\linewidth}}
\toprule
Weighted result & Decision \\
\midrule
\endfirsthead
\toprule
Weighted result & Decision \\
\midrule
\endhead
4.2–5.0 & strong accept / negotiate upward \\
3.5–4.1 & viable if terms are clean \\
3.0–3.4 & negotiate or use as bridge \\
2.5–2.9 & accept only if urgent \\
below 2.5 & reject \\
\bottomrule
\end{longtable}
\end{center}




\subsubsection{Minimum acceptance gates}




Do not accept if any of these fail:




\begin{itemize}
\item no written contract;
\item payment terms unclear;
\item work authorization impossible;
\item unpaid production task required;
\item non-compete blocks future technical art/Unity work broadly;
\item IP terms prevent portfolio use entirely without justification;
\item expected compensation below survival/runway needs.
\end{itemize}




\subsection{16. Counteroffer logic}




\subsubsection{When to counter}




Counter if:




\begin{itemize}
\item the role fit is strong;
\item compensation is below market;
\item contract terms are negotiable;
\item company is responsive;
\item there is no hard legal blocker.
\end{itemize}




\subsubsection{What to counter}




Rank negotiation variables:




\begin{enumerate}
\item base cash compensation;
\item payment schedule;
\item scope/exclusivity;
\item paid trial terms;
\item contract length;
\item notice period;
\item portfolio permission;
\item equipment/software budget;
\item learning budget;
\item relocation or EOR support.
\end{enumerate}




\subsubsection{Counteroffer template}




\begin{quote}
Thank you for the offer. The role is strongly aligned with my Unity and technical-visualization background, especially the work around [specific area]. Based on the scope, contractor structure and market ranges for remote Unity/real-time 3D work, I would be comfortable moving forward at USD [X] gross per month / USD [Y] per hour. I’m also open to discussing structure if there are constraints around salary, such as a paid trial, review after 60–90 days, or a software/equipment budget.
\end{quote}




\subsection{17. Interview preparation calendar}




\subsubsection{Before first applications}




\begin{itemize}
\item Prepare TwinSight 60-second pitch.
\item Prepare TwinSight 2-minute technical explanation.
\item Prepare AI-disclosure answer.
\item Prepare salary-expectation answer.
\item Prepare work-authorization answer.
\item Prepare one C\# code walkthrough.
\item Prepare one optimization story.
\end{itemize}




\subsubsection{Weekly technical practice}




\begin{center}
\scriptsize
\begin{longtable}{>{\raggedright\arraybackslash}p{0.455\linewidth} >{\raggedright\arraybackslash}p{0.455\linewidth}}
\toprule
Area & Frequency \\
\midrule
\endfirsthead
\toprule
Area & Frequency \\
\midrule
\endhead
C\# / Unity fundamentals & 2× week \\
Unity profiling / optimization & 1× week \\
WebGL constraints & 1× week \\
Technical art / rendering & 1× week \\
Portfolio walkthrough & 2× week \\
Mock interview in English & 1× week \\
\bottomrule
\end{longtable}
\end{center}




\subsubsection{Before each interview}




\begin{itemize}
\item Re-read job description.
\item Map role requirements to TwinSight evidence.
\item Prepare three project-specific examples.
\item Check company remote scope.
\item Prepare compensation range.
\item Prepare five questions for interviewer.
\end{itemize}




\subsection{18. Questions to ask employers}




\subsubsection{Technical role clarity}




\begin{itemize}
\item What are the main technical problems this role needs to solve in the first 90 days?
\item Is the role more runtime development, technical art, pipeline tooling or visualization?
\item What platforms are targeted: desktop, web, mobile, XR or embedded?
\item What are the current performance constraints?
\item What is the current asset pipeline?
\end{itemize}




\subsubsection{Remote and contract clarity}




\begin{itemize}
\item Is the role open to Colombia-based contractors?
\item Do you hire through direct B2B, contractor platform or EOR?
\item What countries are eligible for remote work?
\item Are working hours fixed or overlap-based?
\item What is the payment schedule?
\end{itemize}




\subsubsection{Offer safety}




\begin{itemize}
\item What is the governing law of the contract?
\item Are there exclusivity or non-compete clauses?
\item Who owns tools or reusable scripts created during the engagement?
\item Can selected non-confidential work be shown in a portfolio?
\item What is the notice period?
\end{itemize}




\subsection{19. Role-specific proof checklist}




\begin{center}
\scriptsize
\begin{longtable}{>{\raggedright\arraybackslash}p{0.455\linewidth} >{\raggedright\arraybackslash}p{0.455\linewidth}}
\toprule
Role & Must show \\
\midrule
\endfirsthead
\toprule
Role & Must show \\
\midrule
\endhead
Unity Technical Artist & optimization, visual modes, materials, pipeline \\
Real-Time 3D Developer & interaction architecture, C\#, runtime systems \\
Unity WebGL Developer & WebGL build, memory/load constraints \\
Technical Visualization & CAD-to-realtime, metadata, user tasks \\
Digital Twin / Simulation & physical-system reasoning, data limitations \\
XR Developer & XR prototype or strong adjacent proof \\
Tools Developer & automation artifact, editor tooling, Python \\
AI Tools Developer & ARA demo, architecture, reliability limits \\
\bottomrule
\end{longtable}
\end{center}




\subsection{20. Final operational rules}




\begin{itemize}
\item Use TwinSight as the interview anchor.
\item Use ARA only when the role includes tooling, automation or AI workflows.
\item Do not lead with Blender portrait unless applying to 3D-art-adjacent roles.
\item Do not negotiate against Colombian local rates for international contractor roles.
\item Do not accept undefined “remote” terms.
\item Do not complete production-grade unpaid tests.
\item Do not sign broad non-competes without review.
\item Do not overstate EU work authorization before the Portuguese passport or German residence status is real.
\item Do not hide AI assistance if directly asked; frame it as engineering acceleration with ownership.
\end{itemize}




\subsection{21. Source reference list}




Primary source categories used in this section:




\begin{itemize}
\item Internal Revenue Service — guidance on employee vs independent-contractor tax treatment: \url{https://www.irs.gov/businesses/small-businesses-self-employed/independent-contractor-self-employed-or-employee}
\item Lemon.io Unity developer hourly-rate report: \url{https://lemon.io/rate-calculator/unity-developers/}
\item Freelancermap Unity developer role and salary reference: \url{https://www.freelancermap.com/blog/es/que-hace-programador-unity/}
\item Adaface Unity interview question examples: \url{https://www.adaface.com/es/blog/unity-preguntas-entrevista/}
\item Tokio School interview-format reference: \url{https://www.tokioschool.com/noticias/entrevista-tecnica/}
\item Colombia contractor and remote-hiring administrative context from section 03.
\item EU / Portugal / Germany authorization context from section 04.
\item Role-family prioritization from section 02.
\item Portfolio and project-evidence structure from sections 07–10.
\end{itemize}



